\startcomponent ctx-inicio
\product ppchtextut


\startchapter[title={Comando básicos}]

  Empezaremos explicando cómo se utiliza el comando principal para la
  gestión de listas, junto con los argumentos más comunes que modifican
  la salida que producen.
  
  \startsection[title=Comando \color[red]{\type{\chemical}}]
    
    En ConTeXt \MKXL\ ---que es lo mismo que \LUAMETATEX\ ---, el paquete
    \PPCHTEX\ está integrado de forma nativa y el comando más básico es
    \type{\chemical} para escribir reacciones químicas, estructuras y fórmulas
    moleculares.

    Este comando es el componente central del módulo de química \PPCHTEX\ en
    \CONTEXT\, que se usa para representar fórmulas y estructuras químicas.
    Permite la creación de diagramas moleculares complejos, incluyendo anillos y
    enlaces, definiendo enlaces químicos y estructuras atómicas en el entorno
    \type{\startchemical ... \stopchemical}

    La estructura del comando es \type{\chemical[bonds][atoms/molecules]}, donde el
    primer argumento \type{bonds} especifica los enlaces que se dibujarán, por
    ejemplo, \type{SIX} dibuja un anillo de seis enlaces. El segundo argumento
    especifica los átomos que aparecerán, siempre en modo matemático.

    \startsubsection[title=Fórmulas sencillas en modo examen]

      Empezaremos con ejemplos de fórmulas sencillas. Para escribirlas, se pueden
      declarar los subíndices mediante el símbolo \emph{subrayado} del teclado
      \type{_}, o bien mediante \type{\low{subíndice}}.
      Para los superíndices, se puede utilizar el \emph{acento circunflejo}
      \type{^}, o bien \type{\high{superíndice}}.
      
      Veamos algunos ejemplos cuando interesa que unas fórmulas sencillas aparezcan
      en un examen:
      \startitemize[n,packed,joinedup]
      \item Nombra las siguientes especies químicas:
        \startitemize[a, horizontal, five]
          \startitem \chemical{HSO\low{4}\high{-}} \stopitem
          \startitem \chemical{N\low{2}O\low{5}} \stopitem
          \startitem \chemical{CH\low{4}} \stopitem
          \startitem \chemical{Al(CO\low{3})\low{3}} \stopitem
          \startitem \chemical{S\high{2-}} \stopitem
          \startitem \chemical{NH\low{4}NO\low{3}} \stopitem
          \startitem \chemical{C\low{6}H\low{12}O\low{6}} \stopitem
          \startitem \chemical{H\high{-}} \stopitem
          \startitem \chemical{ClO\low{3}\high{-}} \stopitem       
          \startitem \chemical{CH\low{3}-CH\low{2}OH} \stopitem
        \stopitemize
      \item Formula las siguientes especies químicas:
        \startitemize[a, horizontal, two]
          \startitem hidrógenocarbonato de cobalto(II) \stopitem
          \startitem ozono \stopitem
          \startitem benceno \stopitem
          \startitem cloruro de amonio \stopitem  
          \startitem hidrógenofosfato de cadmio \stopitem
          \startitem 1,2-hexadieno \stopitem
          \startitem acetona \stopitem
          \startitem borato de sodio \stopitem
          \startitem ácido sulfhídrico \stopitem
          \startitem nitrógeno \stopitem
        \stopitemize
    \item Nombra las siguientes especies químicas
      \startMycolumns
        \startitemize[a, horizontal, five]
          \startitem \chemical{CO_2} \stopitem
          \startitem \chemical{H_2O} \stopitem
          \startitem \chemical{Cl_2O_7} \stopitem
          \startitem \chemical{Fe_2(SO4)_3} \stopitem
          \startitem \chemical{CH_3-CO-CH_3} \stopitem
          \startitem \chemical{CH_2=CH_2} \stopitem
          \startitem \chemical{HClO_4} \stopitem
          \startitem \chemical{HS^-} \stopitem
          \startitem \chemical{Na_3PO_4} \stopitem
          \startitem \chemical{CH_3NH_2} \stopitem
        \stopitemize
      \stopMycolumns    
    \stopitemize

    \startsubsection[title=Comandos
      \color[red]{\type{\startchemical ... \stopchemical}}]
      
      Este par de comandos se utiliza para dibujar fórmulas más complejas; incluso
      se puede utilizar en su interior múltiples veces el comando \type{\chemical}
      para escribir estructuras mucho más complejas. El siguiente capítulo
      tratará de este tipo de estructuras.
      Por ejemplo, podemos escribir el fenol y el etanol
      \startalignment[center]
      \bTABLE[frame=off, offset=2em, align=center]
      \bTR
      \bTD
      \startchemical[scale=small]
        \chemical [SIX,B,C,+SR6,RZ6][OH]
      \stopchemical
      \eTD
      \bTD
      \startchemical[scale=small]
        \chemical [ONE,Z0357,SB1357,MOV1,Z037,SB137,MOV1,Z01,SB1]
                  [C,H,H,H,C,H,H,O,H]
      \stopchemical  
      \eTD
      \eTR
      \eTABLE
      \stopalignment

  \stopsection

  \startsection[title=\color[red]{\type{\setupchemical}}]
    Ahora presentaremos el comando \type{\setupchemical} dejando de lado el resto
    de parámetros. Con éste podemos influir en la presentación de las fórmulas, a
    menos que en alguna fórmula concreta se especifique otra cosa, como veremos
    en el siguiente párrafo.
    Este comando influye en todas las fórmulas y conviene situarla aparte, por
    ejemplo, en el fichero de entorno general
    % --------------------------------------------------------------------------
    % 1,2,3,4,5,6-hexamentilciclohexano
    % ---------------------------------------------------------------------------
    % Definición de búfer para imagen
    \startbuffer[molecula]
      \setupchemical[axis=on,frame=on,width=5000]
      \startchemical
        \chemical[SIX,B,+SR]
      \stopchemical
    \stopbuffer
    % Definición de búfer para código fuente
    \startbuffer[moleculatxt]
      \setupchemical[axis=on,frame=on,width=5000]
      ...
      \startchemical
        \chemical[SIX,B,+SR]
      \stopchemical
    \stopbuffer
    % .........................................................................
    % Estructura para poner la imagen a la izda. y el código fuente a la dcha.
    % Aquí se utiliza \startlinealignment y \frame
    \startlinealignment[left]
      \framed[frame=off,location=middle]{
        \getbuffer[molecula]}
      \hskip 1.0em % Espacio entre fórmula y código
      \framed[frame=off,location=middle,align=right]
             {\switchtobodyfont[8.0pt]\typebuffer[moleculatxt]}
    \stoplinealignment

    Para que no influya este comando, el argumento del comando se puede añadir
    a \type{\startchemical} de una fórmula concreta, de manera que solo influirá a
    esa fórmula \blank[1ex]
    % --------------------------------------------------------------------------
    % 1,2,3,4,5,6-hexamentilciclohexano
    % ---------------------------------------------------------------------------
    % Definición de buffer
    % Para imagen
    \startbuffer[molecula]
      \startchemical[axis=on,frame=on,width=5000]
        \chemical[SIX,B,+SR,RZ]
      \stopchemical
    \stopbuffer
    % Para código fuente
    \startbuffer[moleculatxt]
      \startchemical[axis=on,frame=on,width=5000]
        \chemical[SIX,B,+SR,-SR,RZ]
      \stopchemical
    \stopbuffer
    % .........................................................................
    % Segunda forma de alinear imagen y código fuente.
    % Aquí se utiliza \setlinealignment y xtable
    \startlinealignment[left] 
    \startembeddedxtable[frame=off, offset=0em, align={right,lohi}]
      \startxrow
        \startxcell
          \getbuffer[molecula]
        \stopxcell
        \startxcell[offset=0.5em]
          {\switchtobodyfont[8pt]\typebuffer[moleculatxt]}
        \stopxcell
          \stopxrow
        \stopembeddedxtable
      \stoplinealignment
    % --------------------------------------------------------------------------

  
\stopchapter
\stopcomponent


%%% Local Variables:
%%% coding: utf-8
%%% mode: context
%%% ConTeXt-Mark-version: "IV"
%%% TeX-engine: default
%%% TeX-master: "../ctx-ppchtextut.tex"
%%% End:
